%% toc: P2 15 15.1,15.2,15.3,15.8
\section{Proofs for decision bounds, reductions and checked certificates}
\label{sec:annex-p02-w02-p02-018-decision-certificates}

Notation is that of \cref{sec:p02-w02-p02-018-decision-certificates}.  For a
valuation \(v\) and a finite formula set \(X\), \(v\entails X\) means
\(v\entails\chi\) for every \(\chi\in X\).  The coincidence lemma
\cref{lem:prop-coincidence} is used in the form: if two valuations agree on
\(\operatorname{Var}(\chi)\), they give \(\chi\) the same value.  A
\emph{step} is one transition of a deterministic multitape machine; a constant
is a number that depends only on the fixed machine being described.

\begin{annexitem}{annex:p02-w02-p02-018-evaluator}{thm:p02-w02-p02-018-evaluator}{Correctness and cost of stack evaluation}
\begin{proof}
Write \(v=v_\tau\).  A variable string \(\texttt{p}\,\nu\,\texttt{;}\) is the
string of a variable exactly when \(\nu\) is a nonempty binary numeral with no
leading \texttt{0} unless it is the single digit \texttt{0}; this is exactly
what item~1 of the method tests.  Let \(\operatorname{red}(\sigma)\) be the
stack obtained by reducing \(\sigma\); reduction only changes entries at the
top, and \(\operatorname{red}(\sigma\cdot\neg\cdot b)=
\operatorname{red}(\sigma\cdot b')\) for a truth value \(b\) with
opposite value \(b'\), because the first reduction step performs
exactly this replacement.

\emph{Formulas are evaluated.}  We prove by structural induction on a formula
\(\chi\): for every stack \(\sigma\) whose top entry is not a truth value,
processing the tokens of \(\chi\) from \(\sigma\) never stops with
\emph{malformed} and ends with \(\operatorname{red}(\sigma\cdot
\widehat v(\chi))\).  If \(\chi\) is a variable or \(\bottom\), the method
pushes \(\widehat v(\chi)\) and reduces.  If \(\chi=\neg\chi'\), it pushes
\(\neg\); the stack \(\sigma\cdot\neg\) has no truth value on top, so by the
induction hypothesis processing \(\chi'\) gives
\(\operatorname{red}(\sigma\cdot\neg\cdot\widehat v(\chi'))=
\operatorname{red}(\sigma\cdot\widehat v(\neg\chi'))\).  If
\(\chi=(\chi_1\circ\chi_2)\), it pushes \((\); by the induction hypothesis
processing \(\chi_1\) gives \(\operatorname{red}(\sigma\cdot(\cdot a)\) with
\(a=\widehat v(\chi_1)\), which equals \(\sigma\cdot(\cdot a\) because the
entry below \(a\) is not \(\neg\).  It then pushes \(\circ\) and, by the
induction hypothesis again, reaches \(\sigma\cdot(\cdot a\cdot\circ\cdot b\)
with \(b=\widehat v(\chi_2)\).  On the final \()\) the four top entries have
the required shape, and the method replaces them by the value of \(a\circ b\),
which is \(\widehat v(\chi)\), and reduces.  Applied to the empty stack, a
formula string ends with the single entry \(\widehat v(\chi)\), so the answer
is \emph{well formed with value \(\widehat v(\chi)\)}.

\emph{Only formulas are accepted.}  Suppose the run on \(w\) does not stop.
We show that after each completed action on a prefix \(u\) of \(w\), with
stack \(e_1\cdots e_h\) from bottom to top, \(u\) splits as
\(u_1\cdots u_h\) where a marker \(e_i\) corresponds to the one-token segment
\(u_i\) equal to it, and a truth value \(e_i\) corresponds to a segment
\(u_i\) that is the string of a formula with \(\widehat v(u_i)=e_i\).  The
empty prefix and empty stack satisfy this.  Pushing a marker adds a
one-token segment.  Item~1 adds the segment \(\texttt{p}\,\nu\,\texttt{;}\),
which is a variable by the first paragraph, with its value under \(v\); item~2
adds \(\bottom\) with value \(\mathbf F\).  A reduction merges the segment
\(\neg\) with the following formula segment \(\chi\) into the formula
\(\neg\chi\), and the new entry is its value.  Item~4 merges the segments
\((\), \(\chi_1\), \(\circ\), \(\chi_2\) and the token \()\) into the formula
\((\chi_1\circ\chi_2)\), and the new entry is its value by the truth clause of
\(\circ\).  If the run ends with one truth value \(b\), the invariant for
\(u=w\) says that \(w\) is a formula string with value \(b\).

Hence \emph{well formed} is answered exactly on formula strings, and then
with the value \(\widehat{v_\tau}(\chi)\).  This proves item~1.

\emph{Cost.}  Use a machine whose read-only input tape holds \(w\) (each token
a fixed group of four cells), whose first work tape holds \(\tau\) as the
string \((n_1,b_1)\cdots(n_r,b_r)\) of \(s\) symbols, and whose second work
tape holds the stack, one cell per entry, with its head on the top entry.
Pushing, popping, and inspecting the top four entries each take a constant
number of steps.  Every reduction or item~4 replacement removes at least one
entry, so their number is at most the number of pushes, which is at most
\(t\).  A lookup for a numeral \(\nu\) scans the table from the left; at an
entry \((n_i,b_i)\) it compares \(n_i\) with \(\nu\) symbol by symbol, moving
both heads forward until a mismatch or the end of one of them, and then
moves the input head back to the start of \(\nu\).  This costs at most
\(2(\operatorname{len}(n_i)+2)\) steps at that entry, where \(\operatorname{len}(n_i)\) is the length of
the numeral, hence at most \(4s\) steps over the whole table, plus \(s\)
steps to return the table head.  There are at most \(t\) lookups.  The total
is at most \(c\,t+t\cdot5s+c\,t\leq c'\,t\,(t+s)\) for constants \(c,c'\),
which proves item~2.
\end{proof}
\end{annexitem}

\begin{annexitem}{annex:p02-w02-p02-018-decision}{thm:p02-w02-p02-018-decision}{Decidability with a running-time bound}
\begin{proof}
\emph{Recognition.}  Decoding fails exactly on bit strings that are not
decodable.  The separator check is a scan of the token string for the shape of
\cref{def:p02-w02-p02-018-input}.  By
\cref{thm:p02-w02-p02-018-evaluator}(1), each segment passes exactly when it
is a formula string.  So Step~1 rejects exactly the malformed inputs, which
belong to none of the four problems.

\emph{Rows.}  Let \(X=\{x_1,\ldots,x_k\}\) be the variables of a well-formed
instance.  Reading \(\mathbf F\) as \(0\) and \(\mathbf T\) as \(1\), the
successor operation of Step~3 is addition of \(1\) to the binary numeral
\(u(x_1)\cdots u(x_k)\) with \(x_k\) least significant, so starting from the
all-\(\mathbf F\) row it visits each of the \(2^k\) assignments
\(X\to\{\mathbf T,\mathbf F\}\) exactly once, in lexicographic order, and
stops after the all-\(\mathbf T\) row.  For \(k=0\) there is one row, the
empty assignment.  Hence Step~3 performs at most \(2^k\) iterations, and the
procedure halts.

\emph{Correctness.}  For a row \(u\) write \(v_u\) for the valuation equal to
\(u\) on \(X\) and \(\mathbf F\) elsewhere; it is the valuation of the table
holding \(u\), so by \cref{thm:p02-w02-p02-018-evaluator}(1) the values
computed in Step~3 are the values \(\widehat{v_u}(\chi)\) of the segments
\(\chi\).  Every valuation \(v\) agrees on \(X\) with \(v_u\) for the row
\(u=v\mathbin{\upharpoonright}X\), and every segment \(\chi\) has
\(\operatorname{Var}(\chi)\subseteq X\), so
\(\widehat v(\chi)=\widehat{v_u}(\chi)\) by the coincidence lemma.  So a
statement of the form ``some valuation gives the segments values with a
property'' holds iff some row does.  Applied to the four problems:
\(\varphi\) is satisfiable iff some row makes it true; \(\varphi\) is valid
iff no row makes it false; \(\{\gamma_1,\ldots,\gamma_r\}\entails\varphi\)
iff no valuation satisfies every \(\gamma_i\) and falsifies \(\varphi\), iff
no row does; and \(\varphi\) and \(\psi\) are equivalent iff no valuation, and
so no row, gives them different values.  Step~3 stops exactly at the first
row with the property, and Step~4 answers when there is none, so the answer is
correct in each case.  The output row \(u\), extended by \(\mathbf F\), is
\(v_u\), which has the property that stopped the search, as claimed.

\emph{Cost.}  Let \(t=m/4\) be the number of tokens.  Decoding and the
separator scan take at most \(c\,m\) steps.  Recognition runs the evaluator on
segments of total length at most \(t\) with the empty table, within
\(c\,t^2\) steps by \cref{thm:p02-w02-p02-018-evaluator}(2).  Step~2 visits at
most \(t\) variable occurrences; for each it compares the numeral with the
list built so far, by the comparison of that proof, within \(c\,t\) steps,
and appends it if it is new; this takes at most \(c\,t^2\) steps.  The table
has at most \(s\leq2t\) symbols, since its numerals have total length at most
\(t\) and each of the \(k\leq t\) entries adds a constant number of symbols.
One row evaluates segments of total length at most \(t\), within
\(c\,t(t+2t)\) steps, keeps a constant amount of information about the values
already seen (for \(\mathsf{CONS}\), whether every premise so far is true; for
\(\mathsf{EQV}\), the first value), and computes the next row within
\(c\,s\) steps.  Summing, the procedure halts within
\(c\,m^2+2^k\,c\,m^2\leq c\,(2^k+1)\,m^2\) steps, for a suitable constant
\(c\).  For a malformed input it halts after recognition, within \(c\,m^2\)
steps, which is the case \(k=0\) of the bound.

\emph{Attainment.}  If \(\varphi\) is valid, no row makes it false, so the
\(\mathsf{TAUT}\) search never stops in Step~3 and visits all \(2^k\) rows.
\end{proof}
\end{annexitem}

\begin{annexitem}{annex:p02-w02-p02-018-limits}{thm:p02-w02-p02-018-limits}{Infinite theories and uncountable alphabets}
\begin{proof}
(1) If \(n\in A\), then \(p_n\in\Gamma_A\), so every valuation satisfying
\(\Gamma_A\) satisfies \(p_n\).  If \(n\notin A\), the valuation \(v\) with
\(v(p_j)=\mathbf T\) iff \(j\in A\) satisfies \(\Gamma_A\) and falsifies
\(p_n\).  Hence \(A=\{n:\Gamma_A\entails p_n\}\), so \(A\) is recovered from
the set \(E_A\) of encodings of consequences of \(\Gamma_A\), and
\(A\mapsto E_A\) is injective.

The family \(\mathcal D\) of languages decided by deterministic multitape
machines is countable.  A machine is given by finitely many tapes, a finite
state set, a finite tape alphabet containing the input symbols and the blank,
and a finite transition table.  Renaming the states and the remaining tape
symbols bijectively by initial segments of \(\N\) does not change the language
decided, and the renamed machine is determined by a finite list of natural
numbers.  The map sending a list \((a_1,\ldots,a_r)\) to the bit string
\(1^{a_1}0\,1^{a_2}0\cdots1^{a_r}0\) is injective, and \(\{0,1\}^*\) is
countable by the enumeration in the proof of (2).  So the machines up to
renaming form a countable set, and \(\mathcal D\) is countable.

Since \(A\mapsto E_A\) is injective, the set \(\{A\subseteq\N:E_A\in
\mathcal D\}\) injects into \(\mathcal D\) and is countable.  By
\cref{thm:p01-cantor} there is no surjection \(\N\to\mathcal P(\N)\), so
\(\mathcal P(\N)\) is uncountable.  Therefore all but countably many
\(A\subseteq\N\) have \(E_A\notin\mathcal D\).

(2) List \(\{0,1\}^*\) by increasing length, and within each length in
lexicographic order; each length class \(\{0,1\}^n\) is finite (it has
\(2^n\) members), so this list assigns every string a unique position in
\(\N\), and the position map is injective.  An injection \(P\to\{0,1\}^*\)
composed with the position map would be an injection \(P\to\N\), making \(P\)
countable.  Hence an uncountable \(P\) has no such injection.
\end{proof}
\end{annexitem}

\begin{annexitem}{annex:p02-w02-p02-018-membership}{thm:p02-w02-p02-018-membership}{\(\mathsf{SAT}\) is in \(\mathsf{NP}\) and \(\mathsf{TAUT}\) is in \(\mathsf{coNP}\)}
\begin{proof}
Let \(n\) be the bit length of \(x\), and take \(q(n)=n\).  For a
well-formed formula instance \(x\) with variables \(x_1,\ldots,x_k\) in order
of first occurrence and a bit string \(y\) of length \(k\), let \(u_y\) be the
row with \(u_y(x_i)=\mathbf T\) iff the \(i\)th bit of \(y\) is \(1\).

Define \(R(x,y)\) to hold iff \(x\) encodes a formula \(\varphi\),
\(\operatorname{len}(y)=k\), and \(\widehat{v_{u_y}}(\varphi)=\mathbf T\).  A machine
deciding \(R\) on \(\langle x,y\rangle\) separates the two parts in one scan
(\cref{def:p02-w02-p02-018-verifier}), runs Steps~1 and~2 of
Method~\ref{def:p02-w02-p02-018-truth-table} on \(x\), compares \(\operatorname{len}(y)\)
with \(k\), copies the bits of \(y\) into the table, and evaluates once.  By
the cost part of the proof of \cref{thm:p02-w02-p02-018-decision} this takes
at most \(c\,(n^2+\operatorname{len}(y)+1)\) steps, a polynomial in the length of
\(\langle x,y\rangle\).  If \(x\in\mathsf{SAT}\), some valuation \(v\)
satisfies \(\varphi\); let \(y\) list \(v(x_1),\ldots,v(x_k)\) as bits.  Then
\(\operatorname{len}(y)=k\leq n\), \(v\) and \(v_{u_y}\) agree on
\(\operatorname{Var}(\varphi)\), so \(R(x,y)\) holds by the coincidence
lemma and \cref{thm:p02-w02-p02-018-evaluator}(1).  Conversely, if
\(R(x,y)\), then \(v_{u_y}\) satisfies \(\varphi\), so \(x\in\mathsf{SAT}\).
Hence \(\mathsf{SAT}\in\mathsf{NP}\) in the sense of
\cref{def:p02-complexity-standing}.

Define \(R'(x,y)\) to hold iff \(x\) is malformed, or \(x\) encodes a formula
\(\varphi\), \(\operatorname{len}(y)=k\) and \(\widehat{v_{u_y}}(\varphi)=\mathbf F\).  It is
decided in the same time, since Step~1 detects malformed inputs.  If
\(x\in\mathsf{TAUT}^{\mathsf c}\), either \(x\) is malformed and
\(R'(x,\varepsilon)\) holds for the empty string \(\varepsilon\), or \(x\)
encodes a formula falsified by some \(v\), and the bits of \(v\) on
\(x_1,\ldots,x_k\) form a \(y\) of length at most \(n\) with \(R'(x,y)\).
Conversely \(R'(x,y)\) implies that \(x\) is malformed or encodes a formula
falsified by \(v_{u_y}\); in both cases \(x\notin\mathsf{TAUT}\).  So
\(\mathsf{TAUT}^{\mathsf c}\in\mathsf{NP}\), and \(\mathsf{TAUT}\), the
complement of an \(\mathsf{NP}\) problem, is in \(\mathsf{coNP}\).
\end{proof}
\end{annexitem}

\begin{annexitem}{annex:p02-w02-p02-018-negation}{thm:p02-w02-p02-018-negation}{The negation reduction}
\begin{proof}
\emph{Computing \(h\).}  Run Step~1 of
Method~\ref{def:p02-w02-p02-018-truth-table} for a single formula, within
\(c\,m^2\) steps.  If \(x\) is well formed, output the four-bit code of
\(\neg\) followed by \(x\); otherwise output the codes of \(\neg\) and
\(\bottom\).  Since \(\neg\varphi\) is a formula whenever \(\varphi\) is, the
output always encodes a formula.

\emph{The equivalences.}  If \(x\) encodes \(\varphi\), then for every
valuation \(v\), \(\widehat v(\neg\varphi)=\mathbf F\) iff
\(\widehat v(\varphi)=\mathbf T\).  So \(\varphi\) is satisfiable iff some
valuation falsifies \(\neg\varphi\), that is iff \(h(x)\notin\mathsf{TAUT}\);
as \(h(x)\) is well formed, this means \(h(x)\in\mathsf{TAUT}^{\mathsf c}\).
Likewise \(\varphi\) is not valid iff some valuation satisfies
\(\neg\varphi\), that is iff \(h(x)\in\mathsf{SAT}\).  If \(x\) is malformed,
then \(x\notin\mathsf{SAT}\) and \(x\in\mathsf{TAUT}^{\mathsf c}\), while
\(h(x)\) encodes \(\neg\bottom\), which every valuation satisfies; so
\(h(x)\in\mathsf{TAUT}\), hence \(h(x)\notin\mathsf{TAUT}^{\mathsf c}\), and
\(h(x)\in\mathsf{SAT}\).  Both displayed equivalences hold in every case.

\emph{Transfer of reductions.}  A machine running within \(p_1(n)\) steps
writes an output of length at most \(p_1(n)\), so composing it with a machine
running within \(p_2(\ell)\) steps on inputs of length \(\ell\) gives a
machine running within \(p_1(n)+p_2(p_1(n))\) steps; polynomial-time maps
are closed under composition.  If \(f\) reduces \(L\) to \(\mathsf{SAT}\),
then \(x\in L\) iff \(f(x)\in\mathsf{SAT}\) iff
\(h(f(x))\in\mathsf{TAUT}^{\mathsf c}\), so \(x\in L^{\mathsf c}\) iff
\(h(f(x))\in\mathsf{TAUT}\), and \(h\circ f\) reduces \(L^{\mathsf c}\) to
\(\mathsf{TAUT}\).  If \(g\) reduces \(L^{\mathsf c}\) to \(\mathsf{TAUT}\),
then \(x\in L\) iff \(g(x)\in\mathsf{TAUT}^{\mathsf c}\) iff
\(h(g(x))\in\mathsf{SAT}\), so \(h\circ g\) reduces \(L\) to
\(\mathsf{SAT}\).
\end{proof}
\end{annexitem}

\begin{annexitem}{annex:p02-w02-p02-018-cook-levin}{thm:p02-w02-p02-018-cook-levin}{Every \(\mathsf{NP}\) problem reduces to \(\mathsf{SAT}\)}
\begin{proof}
\emph{Machine conventions.}  A deterministic machine \(M\) has tapes
\(1,\ldots,\kappa\), each with cells \(0,1,2,\ldots\); a finite state set
\(Q\) with start state \(q_{\mathrm s}\), accepting state \(q_{\mathrm a}\)
and rejecting state \(q_{\mathrm r}\); a finite tape alphabet \(\Sigma_M\)
containing \(0\), \(1\) and the blank \(\sqcup\); and a transition function
assigning to each nonhalting state \(s\) and each tuple
\(\vec a=(a_1,\ldots,a_\kappa)\) of scanned symbols a triple
\(\delta(s,\vec a)=(s',\vec b,\vec d)\) of a next state, symbols to write and
head moves \(d_j\in\{-1,0,1\}\); a move \(-1\) at cell \(0\) leaves the head
at \(0\).  A read-only input tape is the case in which the written symbol on
that tape always equals the scanned one.  Put
\(\delta(s,\vec a)=(s,\vec a,\vec 0)\) for \(s\in\{q_{\mathrm a},
q_{\mathrm r}\}\), so a halted machine stays in its configuration.  On input
\(w\) the run \(c_0,c_1,\ldots\) starts in \(q_{\mathrm s}\) with every head
on cell \(0\), \(w\) in cells \(0,\ldots,\operatorname{len}(w)-1\) of tape \(1\) and blanks
elsewhere, and \(c_{t+1}\) is the successor of \(c_t\) under \(\delta\).  If
\(M\) halts within \(T\) steps, the state of \(c_T\) is the halting state
reached.

\emph{Data.}  Let \(L\in\mathsf{NP}\), witnessed by a polynomial \(q\) and a
predicate \(R\) decided by such a machine \(M\) within \(p(\ell)\) steps on
inputs \(\langle x,y\rangle\) of length \(\ell\)
(\cref{def:p02-w02-p02-018-verifier}).  Replacing \(p\) by a larger
polynomial, assume \(p(\ell)=e\,(\ell+1)^{d}\) with \(e,d\geq1\), so \(p\) is
nondecreasing.  Fix \(x\) of length \(n\), and put
\[
 \ell_0=2n+2,\quad W=q(n),\quad \ell_1=\ell_0+W,\quad T=p(\ell_1),\quad
 N=\ell_1+T.
\]
For every \(y\) with \(\operatorname{len}(y)\leq W\), \(\operatorname{len}(\langle x,y\rangle)\leq\ell_1\),
so \(M\) halts on \(\langle x,y\rangle\) within \(T\) steps.  A head moves at
most one cell per step, so at every time \(t\leq T\) every head is on a cell
at most \(t\leq T<N\).

\emph{Variables.}  For \(0\leq t\leq T\), \(s\in Q\),
\(1\leq j\leq\kappa\), \(0\leq i\leq N\) and \(a\in\Sigma_M\) take distinct
variables \(S_{t,s}\) (state \(s\) at time \(t\)), \(H_{t,j,i}\) (head \(j\) on
cell \(i\) at time \(t\)) and \(C_{t,j,i,a}\) (cell \(i\) of tape \(j\) holds
\(a\) at time \(t\)).  Number \(Q\) and \(\Sigma_M\) from \(0\), let
\(Z=N+\abs Q+\abs{\Sigma_M}+\kappa+3\), and give these variables the indices
\(3(t+Zs)\), \(3(t+Zj+Z^2i)+1\) and \(3(t+Zj+Z^2i+Z^3a)+2\).  Every component
is below \(Z\), so mixed-radix notation and the residue modulo \(3\) make the
indices distinct, and each index has \(O(\log Z)\) binary digits.

\emph{The formula.}  For formulas \(\chi_1,\ldots,\chi_r\) with \(r\geq1\) let
\(\operatorname{One}(\chi_1,\ldots,\chi_r)\) be the right-associated
conjunction of \(\chi_1\lor\cdots\lor\chi_r\) and of the formulas
\(\neg(\chi_a\land\chi_b)\) for \(a<b\); a valuation satisfies it iff it
satisfies exactly one \(\chi_a\).  Let \(\varphi_x\) be the right-associated
conjunction, in a fixed order, of the following formulas, where
\(w=x_1x_1\cdots x_nx_n01\) has length \(\ell_0\).
\begin{enumerate}
\item[(E)] \(\operatorname{One}(S_{t,s}:s\in Q)\) for each \(t\);
      \(\operatorname{One}(H_{t,j,i}:0\leq i\leq N)\) for each \(t,j\); and
      \(\operatorname{One}(C_{t,j,i,a}:a\in\Sigma_M)\) for each \(t,j,i\).
\item[(I)] \(S_{0,q_{\mathrm s}}\); \(H_{0,j,0}\) for each \(j\);
      \(C_{0,1,i,w_{i+1}}\) for \(0\leq i<\ell_0\);
      \(C_{0,1,i,0}\lor C_{0,1,i,1}\lor C_{0,1,i,\sqcup}\) for
      \(\ell_0\leq i<\ell_1\);
      \(C_{0,1,i,\sqcup}\to C_{0,1,i+1,\sqcup}\) for
      \(\ell_0\leq i<\ell_1-1\); \(C_{0,1,i,\sqcup}\) for
      \(\ell_1\leq i\leq N\); and \(C_{0,j,i,\sqcup}\) for \(j\geq2\) and all
      \(i\).
\item[(D)] For each \(t<T\), \(s\in Q\), \(\vec a\in\Sigma_M^\kappa\) and
      \(\vec\imath\in\{0,\ldots,N\}^\kappa\), with
      \(\delta(s,\vec a)=(s',\vec b,\vec d)\) and
      \(i'_j=\max(0,i_j+d_j)\):
      \[
       \Bigl(S_{t,s}\land\textstyle\bigwedge_{j}
         (H_{t,j,i_j}\land C_{t,j,i_j,a_j})\Bigr)\to\Psi,
      \]
      where \(\Psi\) is \(\bottom\) if some \(i'_j>N\), and otherwise
      \(\Psi=S_{t+1,s'}\land\bigwedge_j(H_{t+1,j,i'_j}\land
      C_{t+1,j,i_j,b_j})\).
\item[(F)] \((C_{t,j,i,a}\land\neg H_{t,j,i})\to C_{t+1,j,i,a}\) for each
      \(t<T\), \(j\), \(i\) and \(a\).
\item[(A)] \(S_{T,q_{\mathrm a}}\).
\end{enumerate}

\emph{If \(x\in L\), then \(\varphi_x\) is satisfiable.}  Choose \(y\) with
\(\operatorname{len}(y)\leq W\) and \(R(x,y)\), and let \(c_0,\ldots,c_T\) be the run of
\(M\) on \(\langle x,y\rangle\).  Let \(v\) make \(S_{t,s}\), \(H_{t,j,i}\)
and \(C_{t,j,i,a}\) true exactly when \(c_t\) has state \(s\), head \(j\) on
cell \(i\), and \(a\) in cell \(i\) of tape \(j\).  (E) holds because a
configuration has one state, one position per head and one symbol per cell.
(I) holds because \(\langle x,y\rangle=w\,y\): cells \(\ell_0\) to
\(\ell_1-1\) of tape~\(1\) hold \(y\) followed by blanks, and all later cells
and all other tapes are blank.  In (D), an instance with a true antecedent
lists the state, head positions and scanned symbols of \(c_t\); since the heads
of \(c_{t+1}\) lie on cells at most \(t+1\leq T<N\), every \(i'_j\leq N\), and
\(\Psi\) describes the state, heads and written cells of \(c_{t+1}\), so it
is true.  Instances with a false antecedent are true.  (F) holds because a
cell not under its tape's head is unchanged by a step.  (A) holds because
\(M\) accepts within \(T\) steps and then stays in \(q_{\mathrm a}\).

\emph{If \(\varphi_x\) is satisfiable, then \(x\in L\).}  Let \(v\) satisfy
\(\varphi_x\).  By (E), for each \(t\) there are a unique state
\(\sigma_t\), unique head positions \(\eta_t(j)\leq N\) and unique symbols
\(\gamma_t(j,i)\) made true.  By (I), \(\sigma_0=q_{\mathrm s}\),
\(\eta_0(j)=0\), tape~\(1\) holds \(w\) in cells \(0\) to \(\ell_0-1\), the
cells \(\ell_0\) to \(\ell_1-1\) hold symbols from \(\{0,1,\sqcup\}\) such
that every cell after a blank one is blank, and all other cells are blank.
So these cells hold \(y\,\sqcup^{W-\operatorname{len}(y)}\) for a unique \(y\in\{0,1\}^*\)
with \(\operatorname{len}(y)\leq W\), and \((\sigma_0,\eta_0,\gamma_0)\) is \(c_0\) for the
run on \(\langle x,y\rangle\), restricted to cells \(0,\ldots,N\).
Suppose \((\sigma_t,\eta_t,\gamma_t)\) is \(c_t\) restricted to these cells,
with \(t<T\).  The instance of (D) with \(s=\sigma_t\),
\(i_j=\eta_t(j)\) and \(a_j=\gamma_t(j,\eta_t(j))\) has a true antecedent,
so its \(\Psi\) is true; hence \(\Psi\) is not \(\bottom\), and by the
uniqueness in (E), \(\sigma_{t+1}=s'\), \(\eta_{t+1}(j)=i'_j\) and
\(\gamma_{t+1}(j,i_j)=b_j\).  By (F) and (E), every other cell keeps its
symbol.  Since the scanned symbols of \(c_t\) are the \(a_j\), this is the
successor \(c_{t+1}\) restricted to cells \(0,\ldots,N\).  By induction the
claim holds up to \(T\), and by (A) the state of \(c_T\) is \(q_{\mathrm a}\).
So \(M\) accepts \(\langle x,y\rangle\), \(R(x,y)\) holds with
\(\operatorname{len}(y)\leq q(n)\), and \(x\in L\).

\emph{Polynomial time.}  \(\kappa\), \(\abs Q\) and \(\abs{\Sigma_M}\) are
constants of \(M\), and \(T\), \(N\) and \(Z\) are bounded by a polynomial in
\(n\).  Counting variable occurrences: (D) has
\(T\abs Q\abs{\Sigma_M}^{\kappa}(N+1)^{\kappa}\) formulas with at most
\(3\kappa+1+2\kappa\) occurrences each; the head formulas of (E) have
\(T\kappa\) members with at most \((N+1)^2+(N+1)\) occurrences each; the
remaining formulas of (E), (I), (F) and (A) contribute at most a constant
times \(T(N+1)\).  So \(\varphi_x\) has at most a constant times
\(T(N+1)^{\kappa}+T(N+1)^{2}\) variable occurrences, each written with
\(O(\log Z)\) tokens, and a constant number of other tokens per occurrence;
its length is bounded by a polynomial in \(n\).  A machine writes
\(\varphi_x\) by computing \(n\), \(T\), \(N\) and \(Z\) in binary, and then
running nested loops over the index ranges of (E), (I), (D), (F) and (A),
reading \(\delta\) from its finite table and computing each index by
arithmetic on numerals of \(O(\log Z)\) bits.  Each token written costs a
number of steps polynomial in \(\log Z\), and the four-bit token encoding is
written token by token.  So \(x\mapsto\varphi_x\) is computable in
polynomial time, and it is a polynomial-time many-one reduction of \(L\) to
\(\mathsf{SAT}\).  With \(\mathsf{SAT}\in\mathsf{NP}\)
(\cref{thm:p02-w02-p02-018-membership}), \(\mathsf{SAT}\) is
\(\mathsf{NP}\)-complete.
\end{proof}
\end{annexitem}

\begin{annexitem}{annex:p02-w02-p02-018-taut-conp}{thm:p02-w02-p02-018-taut-conp}{\(\mathsf{TAUT}\) is \(\mathsf{coNP}\)-complete}
\begin{proof}
\(\mathsf{TAUT}\in\mathsf{coNP}\) by \cref{thm:p02-w02-p02-018-membership}.
Let \(L\in\mathsf{coNP}\), so \(L=L'^{\mathsf c}\) for some
\(L'\in\mathsf{NP}\), and \(L^{\mathsf c}=L'\).  By
\cref{thm:p02-w02-p02-018-cook-levin}, \(L^{\mathsf c}\) reduces to
\(\mathsf{SAT}\).  By \cref{thm:p02-w02-p02-018-negation}, applied to the
language \(L^{\mathsf c}\), whose complement is \(L\), the language \(L\)
reduces to \(\mathsf{TAUT}\).
\end{proof}
\end{annexitem}

\begin{annexitem}{annex:p02-w02-p02-018-checks}{thm:p02-w02-p02-018-checks}{Soundness and completeness of the three checks}
\begin{proof}
(1) By \cref{lem:prop-clause-satisfaction}(1), a valuation satisfies a clause
\(C\) iff it makes some literal of \(C\) true, and the value of a literal of
\(C\) depends only on its variable, which lies in \(\operatorname{Var}(F)\).
So every valuation extending the assignment satisfies \(F\) iff every clause
contains a literal true under the assignment, which is what the clause check
tests.  If \(F\) is satisfied by \(v\), the restriction of \(v\) to
\(\operatorname{Var}(F)\) is accepted; if an assignment is accepted, its
extensions satisfy \(F\).

(2) If the log check accepts, the clauses \(C_1,\ldots,C_N\) form a resolution
derivation from \(S\) in the sense of \cref{sec:prop-metatheory-04}: each
\(C_i\) is a member of \(S\) or the resolvent on \(L\) of the earlier
\(C_j\) and \(C_k\).  It contains the empty clause \(C_N\), so \(S\) is not
satisfiable by the soundness theorem of that section.  Conversely, if \(S\)
is finite and not satisfiable, \cref{thm:prop-resolution-complete} gives a
resolution derivation \(D_1,\ldots,D_M\) from \(S\) containing the empty
clause.  Cut it after the first occurrence of the empty clause; a prefix of a
derivation is a derivation.  Justify each member in \(S\) as an input, and
each other member by earlier members and a literal of which it is the
resolvent, which exist by the definition of a derivation.  The resulting log
is accepted.

(3) For \(\chi\in\Gamma\cup\{\varphi\}\), every valuation extending the
assignment agrees with the extension by \(\mathbf F\) on
\(\operatorname{Var}(\chi)\), so gives \(\chi\) the value that the check
computes, by the coincidence lemma.  So the check accepts iff every extension
satisfies \(\Gamma\) and falsifies \(\varphi\).  If such an extension exists,
\(\Gamma\not\entails\varphi\).  If \(\Gamma\not\entails\varphi\), some
valuation \(v\) satisfies \(\Gamma\) and falsifies \(\varphi\), and its
restriction to \(\operatorname{Var}(\Gamma\cup\{\varphi\})\) is accepted.
\end{proof}
\end{annexitem}

\begin{annexitem}{annex:p02-w02-p02-018-extraction}{thm:p02-w02-p02-018-extraction}{Countervaluations from failed searches}
\begin{proof}
(1) The variables listed in Step~2 are exactly those of
\(\Gamma\cup\{\varphi\}\).  By \cref{thm:p02-w02-p02-018-decision} the output
row \(u\), extended by \(\mathbf F\), satisfies every premise and falsifies
\(\varphi\).  Every valuation extending \(u\) agrees with it on
\(\operatorname{Var}(\Gamma\cup\{\varphi\})\), so by
\cref{thm:p02-w02-p02-018-checks}(3) the check accepts \(u\).

(2) By \cref{thm:p12-dpll-correctness}, a returned valuation \(v\) is a total
assignment to \(\operatorname{Var}(F_\Gamma\cup F_\neg)\) satisfying every
clause of \(F_\Gamma\cup F_\neg\).  Let \(v'\) extend it by \(\mathbf F\).
Since satisfaction of a clause depends only on its variables, \(v'\) satisfies
\(F_\Gamma\) and \(F_\neg\), so by hypothesis \(v'\entails\Gamma\) and
\(v'\not\entails\varphi\).  The restriction \(u\) of \(v\) to
\(\operatorname{Var}(\Gamma\cup\{\varphi\})\), which is contained in the
domain of \(v\), agrees with \(v'\) there; every extension of \(u\) therefore
satisfies \(\Gamma\) and falsifies \(\varphi\) by the coincidence lemma, and
the check accepts \(u\) by \cref{thm:p02-w02-p02-018-checks}(3).
\end{proof}
\end{annexitem}

\begin{annexitem}{annex:p02-w02-p02-018-fol-dag}{thm:p02-w02-p02-018-fol-dag}{The annotated refutation is correct}
\begin{proof}
The unification procedure pairs the arguments of \(A\) and \(B\) into
equations \(r\doteq t\), with \(r\) from \(A\) and \(t\) from \(B\).

\emph{First step.}  \(E=C_1\), \(F=C_2\), with
\(\operatorname{Var}(C_1)=\{x\}\) and \(\operatorname{Var}(C_2)=\{y\}\)
disjoint, so the identity renamings rename them apart.  For \(A=Q(x)\) and
\(B=Q(y)\) the equation set is \(\{x\doteq y\}\).  Deletion, decomposition and
orientation do not apply, and elimination applies because
\(x\notin\operatorname{Var}(y)\): it sets \(\delta=\{x\mapsto y\}\), and the
accumulator becomes \(\mathrm{id}\,\delta=\{x\mapsto y\}=\theta_1\) with no
equation left.  Directly, \(Q(x)\theta_1=Q(y)=Q(y)\theta_1\).  If \(\rho\)
unifies \(Q(x)\) and \(Q(y)\), then \(x\rho=y\rho\); so
\((\theta_1\rho)(x)=y\rho=x\rho\) and \((\theta_1\rho)(z)=z\rho\) for
\(z\ne x\), whence \(\rho=\theta_1\rho\) and \(\theta_1\) is a most general
unifier.  The resolvent is
\(\bigl(\{P(x)\}\cup\{R(y)\}\bigr)\theta_1=\{P(y),R(y)\}=C_5\).

\emph{Second step.}  \(E=C_5\), \(F=C_3\); \(C_3\) is ground, so the clauses
are renamed apart.  For \(A=P(y)\) and \(B=P(a)\) the equation set is
\(\{y\doteq a\}\); elimination applies because \(y\notin\operatorname{Var}(a)
=\varnothing\), giving \(\theta_2=\{y\mapsto a\}\).  It unifies, since
\(P(y)\theta_2=P(a)=P(a)\theta_2\).  If \(\rho\) unifies them, then
\(y\rho=a\), so \((\theta_2\rho)(y)=a\rho=a=y\rho\) and \(\theta_2\rho\)
agrees with \(\rho\) elsewhere; thus \(\rho=\theta_2\rho\).  The resolvent is
\(\bigl(\{R(y)\}\cup\varnothing\bigr)\theta_2=\{R(a)\}=C_6\).

\emph{Third step.}  \(E=C_6\) and \(F=C_4\) are ground.  For \(A=B=R(a)\) the
equation set is \(\{a\doteq a\}\); deletion empties it, and the procedure
returns \(\mathrm{id}=\theta_3\).  The identity unifies equal atoms, and every
substitution \(\rho\) equals \(\mathrm{id}\,\rho\).  The resolvent is
\((\varnothing\cup\varnothing)\mathrm{id}=\varnothing=C_7\).

So \(C_1,\ldots,C_7\) is a first-order resolution derivation from
\(\{C_1,C_2,C_3,C_4\}\), whose renamed-apart copies are the clauses
themselves, and it contains the empty clause.  By the soundness theorem for
first-order resolution and factoring in
\cref{sec:fol-cardinality-normalforms-04}, the clause set is not satisfiable.
\end{proof}
\end{annexitem}
